%--------------------------------------------------------------
\documentclass[ja=standard,xelatex]{bxjsarticle}
%--------------------------------------------------------------

%--------------------------------------------------------------
%packages%{
% #region
\usepackage{geometry}
\usepackage{enumitem}
\usepackage{bigints}
\usepackage{changepage}
\usepackage[utf8]{inputenc}
\usepackage{lipsum}
\usepackage[bottom]{footmisc}
\usepackage{subcaption}
\usepackage{zxjatype}
\usepackage[hiragino]{zxjafont} 
\usepackage{tcolorbox}
\usepackage{amsmath}
\usepackage{nccmath}
\usepackage{amsthm}
\usepackage{empheq}
\usepackage{amssymb}
\usepackage{tablefootnote}
\usepackage{bm}
\usepackage{amsfonts}
\usepackage{graphicx}
\usepackage{tocloft}
\usepackage{ascmac}
\usepackage{latexsym}
\usepackage[version=4]{mhchem}
\usepackage{listings, xcolor}
\usepackage{mathtools}
\usepackage[colorlinks=true, linkcolor=blue , citecolor = blue]{hyperref}
\usepackage{array}
\usepackage{url}
\usepackage[english]{babel}
\usepackage[square,numbers]{natbib}
\usepackage{physics}
\usepackage{caption}
\usepackage{ulem}
\usepackage[section]{placeins}
\usepackage{xeCJKfntef}
\usepackage{longtable}
\usepackage{float}
\usepackage{multicol}
\usepackage{mathcomp}
\usepackage{IEEEtrantools}
\usepackage{flafter}
\usepackage{setspace}
% #endregion
%--------------------------------------------------------------%}

%--------------------------------------------------------------
%define%{
% #region

%%skip,space関係
\newcommand{\vs}[1]{\vskip#1\baselineskip}
\newcommand{\parlskip}[1]{\par\setlength{\leftskip}{#1pt}}

%equation環境
\newcommand{\eqn}[1]{\begin{equation*}{#1}\end{equation*}}
\newcommand{\neqn}[1]{\begin{equation}{#1}\end{equation}}
\newcommand{\leqn}[2]{\begin{fleqn}[#1pt]\begin{equation*}{#2}\end{equation*}\end{fleqn}[#1pt]}
\newcommand{\leqnarray}[2]{\begin{fleqn}[#1pt]\begin{eqnarray*}{#2}\end{eqnarray*}\end{fleqn}}
\newcommand{\eqns}[2]{\begin{equation*}\begin{array}{#1}#2\end{array}\end{equation*}}
\newcommand{\neqns}[2]{\begin{IEEEeqnarray}{#1}#2\end{IEEEeqnarray}}
\newcommand{\neqnsequal}{\; =\;}
\newcommand{\neqnarray}[1]{\begin{eqnarray}#1\end{eqnarray}}

%定理・定義・証明等の記述（数学）
\newcommand{\setting}[1]{\vs{0.3}\underline{\textbf{Setting}}\vs{0.3}\parlskip{20}{#1}\parlskip{0}}
\newcommand{\thm}[1]{\vs{0.3}\underline{\textbf{Thm}}\vs{0.3}\parlskip{20}{#1}\parlskip{0}}
\newcommand{\Def}[1]{\vs{0.3}\underline{\textbf{Def}}\vs{0.3}\parlskip{20}{#1}\parlskip{0}}
\newcommand{\rem}[1]{\vs{0.5}\large{\underline{{\textit{remark}}}}\normalsize\vs{0.5}{#1}}
\newcommand{\Proof}[1]{\vs{0.5}\large{\underline{\textit{Proof}}}\normalsize\vs{0.5}{#1}\setcounter{equation}{0}}
\newcommand{\free}[2]{\vs{0.5}\large{\underline{\textit{#1}}}\normalsize\vs{0.5}{#2}\setcounter{equation}{0}}

%tinysections
\newcommand{\tinysection}[3]{\vs{0.5}\large\underline{\textbf{(#1)：#2}}\normalsize\vs{0.5}{#3}}
\newcommand{\tinytinysection}[4]{\vs{0.3}\normalsize{\underline{\textbf{(#1.#2)：#3}}}\vs{0.3}{#4}}
\newcommand{\tinytinytinysection}[5]{\vs{0.2}\normalsize{\underline{(#1.#2.#3)：#4}}\vs{0.2}{#5}}

%微分
\newcommand{\dif}[2]{\dfrac{d#1}{d#2}}
\newcommand{\difop}[1]{\dfrac{d}{d#1}}
\newcommand{\lagdif}[2]{\frac{D#1}{D#2}}

%偏微分
\newcommand{\pdif}[2]{\dfrac{\partial{#1}}{\partial#2}}
\newcommand{\pdifop}[1]{\dfrac{\partial}{\partial#1}}
\newcommand{\pdifl}[1]{\partial_{#1}}
\newcommand{\pdiftwice}[2]{\dfrac{\partial^{2}{#1}}{\partial {#2} ^{2}}}

%積分
\newcommand{\dint}{\displaystyle\int}
\newcommand{\diint}{\displaystyle\iint}
\newcommand{\diiint}{\displaystyle\iiint}

%総和
\newcommand{\sumn}[2]{\sum_{n=#1}^{#2}}
\newcommand{\sumi}[2]{\sum_{i=#1}^{#2}}
\newcommand{\sumj}[2]{\sum_{j=#1}^{#2}}
\newcommand{\sumk}[2]{\sum_{k=#1}^{#2}}
\newcommand{\suml}[2]{\sum_{l=#1}^{#2}}
\newcommand{\summ}[2]{\sum_{m=#1}^{#2}}
\newcommand{\sump}[2]{\sum_{p=#1}^{#2}}
\newcommand{\sumq}[2]{\sum_{q=#1}^{#2}}
\newcommand{\sumfr}[3]{\sum_{#1=#2}^{#3}}
\newcommand{\dsum}{\displaystyle\sum}

%ベクトル
\newcommand{\vect}[1]{\mathbf{#1}}
\newcommand{\vecp}[2]{\mathbf{#1}^{#2}}
\newcommand{\vecn}[1]{\mathbf{#1}^n}
\newcommand{\vecm}[1]{\mathbf{#1}^m}
\newcommand{\pvect}[1]{\begin{pmatrix}#1\end{pmatrix}}

%括弧
\newcommand{\bracketnormal}[1]{\left(#1\right)}
\newcommand{\curlybracket}[1]{\left\{#1\right\}}

%ε-δ論法
\newcommand{\epsilondelta}[4]{\forall\;\epsilon>0,\exists\delta\;\; s.t, \;\; 0<\left| #1-#2 \right|<\delta \Rightarrow \left| #3(#1)-#4\right|<\epsilon}

%その他
\newcommand{\realpower}[1]{\mathbb{R}^{#1}}
\newcommand{\Real}{\mathbb{R}}
\newcommand{\Bf}[1]{\mathbf{#1}}
\newcommand{\limit}[2]{\lim_{#1 \to #2}}
\newcommand{\absdefined}[1]{\left| #1 \right|}
\newcommand{\map}[5]{#1:#2\to#3;\;\;#4\mapsto#5}
\newcommand{\inner}[1]{\left\langle #1 \right\rangle}
\newcommand{\vertline}[2]{\vs{0}\begin{center}\rule{#1\textwidth}{#2pt}\end{center}\vs{0}}
\newcommand{\coord}[2]{\left(#1\;,\;#2\right)}
\newcommand{\scipow}[1]{\times 10^{#1}}
\newcommand{\labeqn}[1]{\label{eqn:#1}}
\newcommand{\labfig}[1]{\label{fig:#1}}
\newcommand{\refeqn}[1]{（\ref{eqn:#1}）}
\newcommand{\reffig}[1]{Figure\ref{fig:#1}}

% #endregion
%--------------------------------------------------------------%}

%--------------------------------------------------------------
%title,author,etc%{
% #region
\title{地球物理数値解析レポート課題（天野教員担当分）}
\author{05-242628三田村彰大}
\date{}
% #endregion
%--------------------------------------------------------------%}

\begin{document}

%--------------------------------------------------------------
% #region
%code_area_settings%{
\lstset{
    basicstyle = {\ttfamily\small}, % 基本的なフォントスタイル
    frame = {tbrl}, % 枠線の枠線。t: top, b: bottom, r: right, l: left
    breaklines = true, % 長い行の改行
    numbers = left, % 行番号の表示。left, right, none
    showspaces = false, % スペースの表示
    showstringspaces = false, % 文字列中のスペースの表示
    showtabs = false, %　タブの表示
    keywordstyle = \color{blue}, % キーワードのスタイル。intやwhileなど
    commentstyle = {\color[HTML]{1AB91A}}, % コメントのスタイル
    identifierstyle = \color{black}, % 識別子のスタイル　関数名や変数名
    stringstyle = \color{brown}, % 文字列のスタイル
    captionpos = t % キャプションの位置 t: 上、b: 下
}
\captionsetup*[lstlisting]{labelformat=empty}
% #endregion
%--------------------------------------------------------------%}

%--------------------------------------------------------------
%document_settings%{
% #region
\setcounter{tocdepth}{3}
\setlength{\parindent}{0pt}
\setlength{\abovedisplayskip}{15pt}
\setlength{\belowdisplayskip}{15pt}
\setlength{\abovedisplayshortskip}{15pt}
\setlength{\belowdisplayshortskip}{15pt}
\setlength{\baselineskip}{15pt}
\newgeometry{top=25truemm,bottom=25truemm,left=20truemm,right=20truemm}
\renewcommand{\arraystretch}{1.5}
\bibliographystyle{abbrvnat}
\setlength{\columnseprule}{1pt}
%asm_settings
\theoremstyle{definition}
\newtheorem{prop}{proposition}[section]
%subfigure_settings
\captionsetup*[subfigure]{position=bottom}
\captionsetup*[figure]{labelfont=bf}
\captionsetup*[table]{labelfont=bf}
%footnote_settings
\setlength{\footnotesep}{0.4\baselineskip} % 脚注どうしの間隔
\makeatletter
\renewcommand\@makefntext[1]{%
  % ---- 脚注本文の段落設定（同一脚注内） ----
  \setlength{\parindent}{0pt}%
  \setlength{\parskip}{0.4\baselineskip}%
  % ---- ぶら下げ＋マーク表示 ----
  \noindent
  \hb@xt@1.8em{\hss\@makefnmark}%
  \hspace{0.4em}%
  \hangindent=2.2em%
  \hangafter=1%
  #1%
}
\renewcommand{\footnoterule}{%
  \kern -3pt
  \hrule width 1.0\linewidth
  \kern 5pt
}
\makeatother
\makeatletter
\renewcommand{\bibsection}{} %"References" 見出しを無効化
\makeatother
\makeatletter
\renewenvironment{proof}[1][\proofname]{%
  \par\pushQED{\qed}\normalfont
  \trivlist
  \item[\hskip\labelsep\bfseries #1\@addpunct{.}]%
  \ignorespaces
}{%
  \popQED\endtrivlist\@endpefalse
}
\makeatother
% #endregion
%--------------------------------------------------------------%}

%--------------------------------------------------------------
%maketitle_makehedding_tableofcontents%{
% #region
\maketitle



\begin{description}[
  labelwidth=10.5cm,
  labelsep=0.5cm,
  leftmargin=5.8cm,
  style=nextline
]
  \item[数値計算結果の可視化動画置き場：] \href{https://shota-mitamura.com/works/geo_analysiskadai1/kadai/}{https://shota-mitamura.com/works/geo\_analysiskadai1/kadai/}
  \item[コード置き場：] \href{https://github.com/ramutami/numeric_geophysics_report1}{https://github.com/ramutami/numeric\_geophysics\_report1}
\end{description}


\vs{5}
\tableofcontents
\vs{3}
※解いた問題：課題A、課題B、課題D-1、課題D-2
\newpage
% #endregion
%--------------------------------------------------------------%}

%start_document
\section{課題A}
\subsection*{全体の感想}
　それぞれのスキームがどのような発想から生まれるかといったことが丁寧に解説されたため、非常に分かりやすかった。また課題で数値計算スキームを可視化したのがとても楽しかった。
\subsection*{線形問題のパートについて}
　スキームの比較が非常に面白かった。特に数値拡散まわり（Lax-Wendroffや人工粘性など）の話は明快でよく理解することができたように思う。ただ、保存スキームのあたりは最初は少し理解に苦しんだ。（有限体積法を知らなかった為だと思われる。）
\subsection*{非線形問題のパートについて}
　Riemann問題の話は非常に面白かった。ただしその前段階の弱解やエントロピー条件などの話は難しいと感じた。元の微分方程式が滑らかな関数の範囲でとけない時に積分系を満たす解から物理的に適切な解を持ってくるという発想は理屈としては納得できるが、厳密になぜそのやり方で正しい解を拾ってこれるのかが少しモヤモヤした。
　\vs{0}
　ただ、euler方程式周りの話は非常にわかりやすかった。（元々馴染みのある方程式だったからかもしれないが…）できればレポートのeuler方程式の問題も解きたかったが、時間が足りなくて断念してしまったのが悔やまれる。
\section{課題B}
\begin{spacing}{2.0}
（B-1）$u\bracketnormal{x,t} = u_0\bracketnormal{x-tc}$　（B-2）$\pdif{u}{t} = \dfrac{u_i^{n+1}-u_i^n}{\Delta t}$　（B-3）$\pdif{u}{x}=\dfrac{u_{i+1}^n-u_{i-1}^n}{2\Delta x}$　
\vs{0}
（B-4）$u_i^{n+1} = u_i^n - \dfrac{c\Delta t}{2\delta x}\bracketnormal{u_i^{n+1}-u_{i-1}^n} $　（B-5）無条件不安定　
\vs{0}（B-6）$0\leq \nu \leq 1$なら安定、ただし$\nu = \dfrac{c\Delta t}{\Delta x}$　（B-7）$-1\leq \nu \leq 0$なら安定、ただし$\nu = \dfrac{c\Delta t}{\Delta x}$　
\vs{0}
\end{spacing}
（B-8）数値計算スキームが安定となるためには、物理的な情報の伝播がすでに済んでいる位置での情報を使って求めたい位置での時間ステップ更新を行う必要がある。
\begin{spacing}{2.0}
（B-9）位相誤差　（B-10）位相速度のズレが大きくなる　（B-11）$\dfrac{1}{2}c^2 \pdiftwice{u}{x}\Delta t$　（B-12）逆拡散　
\vs{0}
（B-13）数値拡散の項（あるいは、$\curlybracket{\dfrac{1}{2}\nu ^2\bracketnormal{u_{i+1}^n - 2u_{i}^n+u_{i-1}^n}}$の項）　（B-14）数値振動　
\vs{0}
（B-15）Godunovの定理　（B-16）人工粘性　（B-17）積分　（B-18）保存則が成立　（B-19）特性量　
\vs{0}
（B-20）$\dif{x}{t}=\lambda\bracketnormal{x}$　（B-21）線形移流方程式が$\pdif{u}{t}+c\pdif{u}{x}$と書き表される時、移流速度$c$　（B-22）物理量$u$　
\vs{0}
（B-23）２つの特性量の値を持つことが必要　（B-24）衝撃波　（B-25）粘性（数値拡散）　（B-26）弱解　
\vs{0}
（B-27）エントロピー条件　（B-28）左（負）向きの音波　（B-29）エントロピー波　（B-30）右（正）向きの音波
\end{spacing}

\newpage
\section{課題D-1}
　$x>0$向きの移流速度をもつ線形移流方程式を考えることにする。すなわち次の方程式を考える。
\neqn{
  \pdif{u}{t} + c\pdif{u}{x} = 0 \quad \bracketnormal{c>0}
}
ただし$c$は計算に際して適宜定めるものとする。
\vs{0}
　また、今回は一次精度風上差分とLax-Wendroffを用いて上式を解くことを考える。具体的には、空間方向に幅$\Delta x$で$i=1,2,\ldots$と離散化し、時間方向に幅$\Delta t$で$n=0,1,\ldots$と離散化する時、一次精度風上差分及びLax-Wendroffは（$c>0$よりFTBSとなるので）次のようなスキームとなる。
\neqns{ll}{
  \text{一次精度風上差分（FTBS）:}\quad & u_i^{n+1} = u_i^n - \dfrac{c\Delta t}{\Delta x}\bracketnormal{u_i^n - u_{i-1}^{n}}\\[5pt]
  \text{Lax-Wendroff:}\quad & u_i^{n+1} = u_i^n - \dfrac{c\Delta t}{2\Delta x}\bracketnormal{u_{i+1}^n - u_{i-1}^{n}} + \dfrac{1}{2}\bracketnormal{\dfrac{c\Delta t}{\Delta x}}^2 \bracketnormal{u_{i+1}^n - 2u_i^n + u_{i-1}^n}
}
　実際の計算においては計算領域を$0\leq x \leq 1$とした上で周期境界を用い、初期条件は$u_0(x) =\sin \bracketnormal{2\pi x}$と定める。また、$c=0.5$と定め、適当に定めた格子幅$\Delta x$に対しCFL条件を満たすように（すなわち$c\Delta t/\Delta x <1 $となるように）$\Delta t$を定めた。\footnote{
  具体的には、$\Delta x=0.01$、$c\Delta t/\Delta x = 1/3$とした。
}
また計算はFortranを用いて行った。\footnote{
  Fortranが一番慣れているので...
}
そのようにして計算して得られた結果を\reffig{1}左図に示す。
\begin{figure}[!htbp]
    \centering
    \begin{subfigure}{0.45\textwidth}
        \centering
        \includegraphics[width = \linewidth]{./image_tex/image000150.png}
    \end{subfigure}
    \hspace{0.05\linewidth}
    \begin{subfigure}{0.45\textwidth}
        \centering
        \includegraphics[width=\linewidth]{./image_tex/image_box003500.png}
    \end{subfigure}
    \caption{一次精度風上差分及びLax-Wendroffを用いた移流の数値計算結果。左は正弦波を初期条件に$t=1.0$まで計算した結果。右は矩形波を初期条件に$t=7.0$まで計算した結果。（\href{https://shota-mitamura.com/works/geo_analysiskadai1/kadai/}{動画ページ}ではそれぞれ図１及び図２に対応している。）}
    \label{fig:1}
\end{figure}

\vs{0}
　今\reffig{1}左図を見ると、1次元風上差分（図中青線）では数値拡散によって波のピークが低くなってしまっているのに対し、Lax-Wendroff（図中オレンジ線）ではそのような減少が生じていない。これをより詳しく見るために、矩形波を初期条件として計算した結果を示したのが\reffig{1}右図である。ただしここでは初期条件を
\neqn{
  u_0\bracketnormal{x} = \left\{\begin{aligned}
    -1\quad (x<0.5)\\
    1 \quad (x>0.5)
  \end{aligned}\right.
}
とした。\reffig{1}右図を見ると、一次精度風上差分の数値拡散が非常に良く見て取れる。また、Lax-Wendroffは時間が経っても波の強さを維持することができているが、同時に数値振動が生じてしまっている。
\FloatBarrier
\vs{1}
　次に、$\Delta x$の値を様々に変化させて誤差
\neqn{
  \epsilon = \sqrt{\dint\absdefined{u\bracketnormal{x}-u_\text{analytic}\bracketnormal{x}}^2 dx}\approx\sqrt{\dsum_i \absdefined{u_i^n-u_\text{analytic}\bracketnormal{x_i}}^2 \Delta x}\labeqn{2}
}
がどのように変化するかを調べる。そのためにまず、それぞれのスキームにおける誤差の理論式を導出する。
\vs{1}
　まず一次精度風上差分について考える。安定性解析をすると、複素増幅率は次のように求まる。
\neqn{
  g = \dfrac{\tilde{u}^{n+1}}{\tilde{u}^n} = 1-\nu \bracketnormal{1-e^{-j\theta}}
}
ただし$\nu = c\Delta t/\Delta x$,$\theta = k\Delta x$であり、$k$はフーリエ展開における波数、$j$は虚数単位である。
\vs{0}
　この時、初期条件の正弦波は$k=2\pi$の波であり、その複素振幅は$1$である。よって十分時間が経ったのちも$k=2\pi$であると考えれば、この波の複素振幅は$\tilde{u}^n = 1\times g^n$となるはずである。よって$t=T=n\Delta t$での波は$u_\text{upwind}^n =\Im{g^n e^{jk\Delta x}}$となる。（ただし$k=2\pi$である。）また数値解は$u_\text{analytic}=\sin\bracketnormal{2\pi\bracketnormal{x-cT}}$であるが、これをフーリエ展開の形で書くと$u_\text{analytic} = \Im{e^{j2\pi (x-cT)}}$である。
\vs{0}
　よって\refeqn{2}中に現れる値は次のように計算される。
\neqns{lcl}{
  \absdefined{u_i^n-u_\text{analytic}\bracketnormal{x_i}}^2 &\; =\;& \absdefined{\Im{g^n e^{jk\Delta x}}-\Im{e^{j2\pi (x-cT)}}}^2\\
  &\; =\;& \absdefined{\Im{e^{j2\pi x}\bracketnormal{{e^{-j2\pi cT}-g^n} }}}^2
}
よって$e^{-j2\pi cT}-g^n=A=\absdefined{A}e^{i\phi}$とおけば次を得る。
\neqn{
  \absdefined{u_i^n-u_\text{analytic}\bracketnormal{x_i}}^2 = \absdefined{A}^2 \sin^2\bracketnormal{\phi + 2\pi x}
}
よって、\refeqn{2}中の積分にこれを代入すれば、誤差は次のように計算される。
\neqn{
  \epsilon_\text{upwind}  = \sqrt{\dint{\absdefined{u_i^n-u_\text{analytic}\bracketnormal{x_i}}^2}}  =  \absdefined{A}\curlybracket{{\dint_0^1\sin^2\bracketnormal{\phi + 2\pi x}dx}}^{1/2} = \dfrac{\absdefined{A}}{\sqrt{2}}\labeqn{3}
}
　特に今、$e^{-j2\pi cT} = \bracketnormal{e^{-j2\pi \nu \Delta x}}^n $であるから、これを$E^n$と書くことにすれば、\refeqn{3}式中に現れる$A$を次のように計算することができる。
\neqns{lcl}{
  A &\; = \; & E^n-g^n\\
  &\; = \; & \bracketnormal{g-E}\bracketnormal{g^{n-1}+g^{n-2}E + \cdots +E^{n-1}}\\
  &\; \vdots \; &  \\
  &\; = \; & -\dfrac{1}{2}\bracketnormal{1-\nu}k^2 c\Delta x E^{n-1} + \mathcal{O}\bracketnormal{\Delta x^2}
}

\footnote{
  ここの計算においては、$g$及び$E$を
  \neqns{l}{
    g = 1-\nu\bracketnormal{1-e^{-jk\Delta x}}= 1-\nu j2\pi \Delta x - \dfrac{\nu}{2} k^2 \Delta ^2 + \mathcal{O}\bracketnormal{\Delta x ^3} \\
    E = e^{-j2\pi \nu \Delta x}= 1-\nu j2\pi \Delta x - \dfrac{\nu^2}{2} k^2 \Delta ^2 + \mathcal{O}\bracketnormal{\Delta x ^3}
  }
  のようにテイラー展開し、そこから従う
  \neqns{l}{
    g-E = -\dfrac{\nu\bracketnormal{1-\nu}}{2}\bracketnormal{k\Delta x}^2+\mathcal{O}\bracketnormal{\Delta x^3}\\
    g^{n-1}+g^{n-2}E + \cdots +E^{n-1} = nE^{n-1} + \mathcal{O}\bracketnormal{n^2\Delta x^2} = nE^{n-1} + \mathcal{O}\bracketnormal{\bracketnormal{cT/\nu}^2} \sim nE^{n-1} + \mathcal{O}\bracketnormal{1}
  }
  などを用いている。（行間がそこそこ長いので割愛するが、間違っていそうで怖い…）
}
よって$\absdefined{E}=1$などを用いれば、$\epsilon_\text{upwind} = \absdefined{A}/\sqrt{2}$は次のように求まる。
\neqn{
  \epsilon_\text{upwind} = \dfrac{c}{2\sqrt{2}}\bracketnormal{1-\nu}k^2\Delta x + \mathcal{O}\bracketnormal{\Delta x^2}
}
よって、$\Delta x \ll 1$で
$
  \epsilon_\text{upwind} \propto \Delta x
$
が成立する。
\vs{0}
　以上の考察から最終的な$\epsilon_\text{upwind}$と$\Delta x$の両対数プロットは傾き$1$の直線になると期待されるが、$\Delta x \sim 1$においては直線から離れることも予想される。
\vs{1}
　次にLax-Wendroffの場合について考える。この時、安定性解析をすると複素増幅率は次のように求まる。
\neqn{
  g = \dfrac{\tilde{u}^{n+1}}{\tilde{u}^n} = 1-j\nu \sin\bracketnormal{k\Delta x}+ \nu ^2 \bracketnormal{\cos\bracketnormal{k\Delta x}-1}\labeqn{18}
}
この場合においても、$t=T$での解析解と数値解の誤差は
\neqn{
  \epsilon_\text{Lax-Wendroff} = \dfrac{\absdefined{A}}{\sqrt{2}} \quad \bracketnormal{\text{with:}\; A = E^n-g^n\;,\; E = e^{-j2\pi \nu \Delta x}}
}
のように計算されることは変わらないので、上式の$g$に\refeqn{18}を代入すればLax-Wendroffの場合の誤差を得ることができる。$A$については先と同様に計算することができ、
\neqn{
  A  =  E^n - g^n =\cdots = \dfrac{j}{6}\bracketnormal{\nu - \nu^3}\dfrac{Tk c}{\nu}\bracketnormal{k\Delta x}^2 E^{n-1} + \mathcal{O}\bracketnormal{\Delta x^3}
}
となる。\footnote{
  ここの計算においても、（$k\Delta x=\theta$とおいた上での）テイラー展開
  \neqns{l}{
    g = 1-j\nu \sin\bracketnormal{k\Delta x}+ \nu ^2 \bracketnormal{\cos\bracketnormal{k\Delta x}-1} = 1-j\nu \theta -\dfrac{\nu^2}{2} \theta ^2 + \dfrac{j\nu}{6}\theta^3 + \mathcal{O}\bracketnormal{\theta^4} \\
    E = e^{-j2\pi \nu \Delta x} = e^{-j \nu \theta} = 1-j\nu\theta - \dfrac{\nu^2}{2}\theta^2 + \dfrac{j\nu^3}{6}\theta^3 + \mathcal{O}\bracketnormal{\theta^4}
  }
  を用いることで得られる
  \neqns{l}{
    g-E = \dfrac{j}{6}\bracketnormal{\nu - \nu^3}\theta^3 + \mathcal{O}\bracketnormal{\theta^4}\\
    g^{n-1}+g^{n-2}E + \cdots +E^{n-1} = nE^{n-1} + \mathcal{O}\bracketnormal{1}
  }
  と言う関係を用いている。
}
よって先と同様にして、$\epsilon_\text{Lax-Wendroff}$として次を得る。
\neqn{
  \epsilon_\text{Lax-Wendroff} = \dfrac{1}{6\sqrt{2}}\bracketnormal{\nu - \nu ^3}\dfrac{Tk^3c}{\nu}\Delta x^2 + \mathcal{O}\bracketnormal{\Delta x^3}
}
よって、$\Delta x \ll 1$で$\epsilon_\text{Lax-Wendroff}\propto \Delta x^2$が成立する。
\vs{0}
　以上の考察から、最終的な$\epsilon_\text{Lax-Wendroff}$と$\Delta x$の両対数プロットは傾き$2$の直線になると期待されるが、やはり$\Delta x\sim 1$においては直線から離れることが予想される。
\vs{1}
　ここまでの考察を踏まえた上で、実際に$\Delta x$を様々に変えながら$\epsilon$を計算し、プロットした図が\reffig{4}。

\begin{figure}[!htbp]
  \centering
  \includegraphics[width=0.5\textwidth]{./image_tex/epsilon.png}
  \caption{解析解と数値解の誤差$\epsilon$の$\Delta x$依存性。青が一次精度風上差分、オレンジがLax-Wendroff。青の波線は$\epsilon \propto \Delta x$の直線を、オレンジの破線は$\epsilon \propto \Delta x^2$の直線を表している。}
  \label{fig:4}
\end{figure}
\vs{0}
　\reffig{4}を見ると、一次精度風上差分は$\epsilon\propto\Delta x$となっており、Lax-Wendroffは$\epsilon\propto\Delta x^2$となっていることがわかる。（それぞれの色の破線と傾きが一致していることからわかる。）よって題意は示せた。また、$\Delta x\sim 1$近傍で$\epsilon$が$\Delta x$や$\Delta x^2$に比例しなくなるのも予想通りである。（傾きが若干下がってしまっている。）先に導いた誤差の理論式を考えれば、これは$\Delta x$が大きい時は誤差の高次の項が無視できなくなることが原因であるとわかる。
\FloatBarrier
\subsection*{おまけ：その他のスキーム}
　今回、線形移流方程式でにおいて一次精度風上差分とLax-Wendroffの挙動を観察したが、そのついでにFTCSやKawahara-kuwaharaスキーム、QUICKESTスキームなどでも数値解の挙動を確認した。\footnote{
  Kawahara-Kuwaharaスキームは工学部の災害シミュレーションの授業でちょうど扱ったため、採用した。またQUICKESTスキームは名前がカッコよかったので採用した。（あとすごいらしいと聞いたので。）
}
\footnote{
  Kawamura-Kuwaharaスキームは\cite{Kawakuwa}
  を、QUICKESTスキームは\cite{QUICKEST}を参考にした。
}

\neqns{ll}{
  \text{FTCS:}\qquad & u_i^{n+1}=u_i^n-\frac{\nu}{2}\left(u_{i+1}^n-u_{i-1}^n\right)\\[10pt]
  \text{Kawahara-Kuwahara:} \quad& u_i^{n+1}=u_i^n-\dfrac{\nu}{12}\left(-u_{i+2}^n+8u_{i+1}^n-8u_{i-1}^n+u_{i-2}^n\right)\nonumber \\
  & \qquad \qquad \quad  -\dfrac{\nu}{4}\left(u_{i+2}^n-4u_{i+1}^n+6u_i^n-4u_{i-1}^n+u_{i-2}^n\right)\\[10pt]
  \text{QUICKEST:} & u_i^{n+1}=\frac{\nu(\nu^2-1)}{6}u_{i-2}^n-\frac{\nu(\nu^2-\nu-2)}{2}u_{i-1}^n \nonumber \\
  & \qquad \qquad \quad  +\left\{1+\frac{\nu}{2}\left(\nu^2-2\nu-1\right)\right\}u_i^n-\frac{\nu(\nu^2-3\nu+2)}{6}u_{i+1}^n 
}

その結果を示したのが\reffig{5}である。
\begin{figure}[!htbp]
    \centering
    \begin{subfigure}{0.31\textwidth}
        \centering
        \includegraphics[width = \linewidth]{./image_tex/image_box_FTCS003500.png}
    \end{subfigure}
    \begin{subfigure}{0.31\textwidth}
        \centering
        \includegraphics[width=\linewidth]{./image_tex/image_box_KAWAKUWA003500.png}
    \end{subfigure}
    \begin{subfigure}{0.31\textwidth}
        \centering
        \includegraphics[width=\linewidth]{./image_tex/image_box_QUICKEST003500.png}
    \end{subfigure}
    \caption{様々なスキームにおける矩形波の移流の計算結果。左からFTCSスキーム、Kawahara-Kuwaharaスキーム、QUICKESTスキーム。}
    \label{fig:5}
\end{figure}

\vs{0}
　\reffig{5}左図中の黄色線は、矩形波を初期条件にFTCSスキームで時間発展させた時の数値解を示している。この図を見ると解が激しく振動していることがわかる。また、この後も時間発展を続けると、解がやがて発散することがわかり、FTCSスキームが無条件不安定となっていることが見て取れる。（\href{https://shota-mitamura.com/works/geo_analysiskadai1/kadai/}{動画ページ}の図３を参照。）
\vs{0}
　また、\reffig{5}右図中の緑色線は、同様の初期条件をKawamura-Kuwaharaスキームで発展させた時の様子を示している。この図を見ると、FTCSのような不安定さや一次精度風上差分のような数値拡散がなく、尚且つLax-Wendroffのような激しい数値振動も生じていないことがわかる。しかし、実際にこの結果を得るためには格子間隔をFTCSに比べ優に小さくする必要があり、FTCSと同様の計算量では解は不安定となってしまった。（\href{https://shota-mitamura.com/works/geo_analysiskadai1/kadai/}{動画ページ}の図４を参照。）
\vs{0}
　最後に、\reffig{5}右図中の薄紫色線は、同様の初期条件をQUICKESTスキームで発展させた時の様子を示している。この図を見ると、QUICKESTスキームにおいてもまた、Kawamura-Kuwaharaスキーム同様不安定さや数値拡散・数値振動の問題が全て抑えられていることがわかる。その上で、QUICKESTスキームで安定な解を得るためにはKawamura-Kuwaharaスキームほどの小さな格子間隔を必要としなかった。この点から、QUICKESTスキームが計算量という観点においても優秀であることを体感することができた。（\href{https://shota-mitamura.com/works/geo_analysiskadai1/kadai/}{動画ページ}では図５に対応。）
\FloatBarrier
\vs{1}
　最後に、これまで出てきた全てのスキームの比較する。今\reffig{6}に、全てのスキームによる時間発展の様子を示している。またTable\ref{tab:1}にはそれぞれのスキームの特徴をまとめている。今回の実験で、直感的には
\neqn{
  \text{一次精度風上差分}<\text{FTCS}<\text{Lax-Wendroff}<\text{Kawamura-Kuwahara}<\text{QUICKEST}
}
の順にスキームが「高級」になっていっていることがわかった。（この順に、一つ前にあるスキームの問題を解決していっている。（Table\ref{tab:1}の備考参照。））

\begin{figure}[!htbp]
  \centering
  \includegraphics[width=0.7\textwidth]{./image_tex/image_box_all003500.png}
  \caption{線形移流方程式に従う矩形波の時間発展の様子。黒が解析解、水色が一次精度風上差分スキーム、オレンジがLax-Wendroffスキーム、黄色がFTCSスキーム、緑がKawamura-Kuwaharaスキーム、薄紫がQUICKESTスキームに対応している。（\href{https://shota-mitamura.com/works/geo_analysiskadai1/kadai/}{動画ページ}では図６に対応。）}
  \label{fig:6}
\end{figure}
\vs{0}

\begin{table}[!htbp]
    \centering
    \begin{tabular}{|c||c|c|c|c|c|p{4cm}|}
        \hline
        
        スキーム& 数値拡散 & 安定性 & 数値振動 & 計算量 & 空間精度 & 備考\\
        \hline\hline
        一次精度風上差分 & $\times$ &  &  & & $\times$& 1次精度しかない\\
        \hline
        FTCS &  & $\times$ & $\times$ &  & $\circ$& 2次精度あるが、不安定\\
        \hline
        Lax-Wendroff & $\circ$ & $\circ$ & $\times$ &  & $\circ$& 2次精度かつ安定だが、数値振動\\
        \hline
        Kawamura-Kuwahara & $\circ$ & $\circ$ & $\circ$ & $\times$ & $\circ\circ$&3次精度\tablefootnote{\cite{Kawakuwa}によれば、3次精度風上差分となっているらしい}かつ安定かつ数値振動しないが、計算量$\times$\\
        \hline
        QUICKEST & $\circ$ & $\circ$ & $\circ$ & $\circ$ & $\circ\circ$&3次精度\tablefootnote{\cite{QUICKESTorder}によればQUICKESTは3次らしい…？（要確認）}かつ安定かつ数値振動しないかつ計算量$\circ$\\
        \hline
    \end{tabular}
    \caption{スキームの性質の比較。$\circ$は性質が良いことを表し、$\times$は性質が悪いことを表す。}
    \label{tab:1}
\end{table}



\FloatBarrier

\newpage
\section{課題D-2}
　非粘性Burgers方程式
\neqn{
  \pdif{u}{t} + u\pdif{u}{x} = 0
}
を与えられた初期条件
のもとで解くことを考える。ただし、問題では計算領域を$-1\leq x\leq 1$としているが、より詳しく波の性質を観察するため計算領域は$-2\leq x \leq 2$と取ることにする\footnote{
  計算領域が$-1\leq x\leq 1$だと、周期境界を回り込んだ波がすぐに自分の波の後端に追いついてしまい、重ね合わせが生じてしまう。そこで、できるだけ長く波の挙動を観察するために計算量域を長めに取ることにした。
}。すなわち初期条件を
\neqn{
  u\bracketnormal{x} = \left\{\begin{aligned}
    0 && -1<x<-1/3 \\
    1 && -1/3<x<1/3\\
    0 && 1/3 < x < 1
  \end{aligned}\right.
  \qquad \leadsto\qquad
  u\bracketnormal{x} = \left\{\begin{aligned}
    0 && -2<x<-1/3 \\
    1 && -1/3<x<1/3\\
    0 && 1/3 < x < 2
  \end{aligned}\right.
}
と変更することにする\footnote{領域を広げただけで本質的には何も変えていない。}。
\subsection*{解析解}
　まず、この場合の解析解を考えることにする。そのためにはまず積分形を満たす弱解を構成し、その解がエントロピー条件を満たしているかどうかを確認すれば良い。
\vs{0}
　その上で、今回は初期条件が矩形波であり、初め$x=1/3$にある不連続面を前面、$x=-1/3$にある不連続面を後面と呼ぶことにすると、前面においては$U_L>U_R$のRiemann問題を考えればよく、後面においては$U_L<U_R$のRiemann問題を考えれば良い。（\reffig{7}参照。）
\begin{figure}[!htbp]
  \centering
  \includegraphics[width=0.7\textwidth]{./image_tex/riemannseperation.png}
  \caption{課題D-2の問題の読み替え。二つの不連続面はそれぞれ$U_L>U_R$及び$U_L<U_R$のRiemann問題に対応していると考えることができ、よって前面は衝撃波、後面は膨張波となるような解が解析解となることが期待される。}
  \label{fig:7}
\end{figure}
\vs{0}
　よって、このようなRiemann問題の解析解を考えれば、今回の場合においては前面が衝撃波で後面が膨張波となるような波が解析解となることが期待される。すなわち今、\reffig{7}のように後面のRiemann問題を$\bracketnormal{U_{L_1},U_{R_1}}$、前面のRiemann問題を$\bracketnormal{U_{L_2},U_{R_2}}$のRiemann問題として定式化することを考えれば、$U_{L_1} = U_{R_2} = 0$かつ$U_{R_1} = U_{L_2} = 1$であり、これらに対し解析解は次のようになっているべきである。
\neqn{
  u(x) = \left\{
  \begin{aligned}
    \;&U_{L_1}  && \Big(\; x<-\dfrac{1}{3} + U_{L_1} t\;\Big)\\
    &\dfrac{x-(-1/3)}{t}  \qquad && \Big(\;-\dfrac{1}{3} + U_{L_1} t<x<-\dfrac{1}{3} + U_{R_1} t\;\Big)\\
    &U_{R_1} && \Big(\; -\dfrac{1}{3} + U_{R_1} t< x \;\Big)\\
    &U_{L_2} && \Big(\;  x< +\dfrac{1}{3} + St\;\Big)\\
    &U_{R_2} && \Big(\;  +\dfrac{1}{3} + St <x \;\Big)
  \end{aligned}\right. \labeqn{34}
}
\FloatBarrier
ただし衝撃は速度$S$はBurgers方程式における衝撃波速度の条件$S = \bracketnormal{U_{L_2}+U_{R_2}}/2$を満たさねばならない。よってこれらの式に$U_{L_1}$などの値を代入していけば、予想される解析解は\refeqn{35}式及び\reffig{8}で表されるようなものとなる。
\neqn{
  u(x) = \left\{
  \begin{aligned}
    \;&0  && \Big(\; x<-\dfrac{1}{3}\;\Big)\\
    &\dfrac{x+1/3}{t}  \qquad && \Big(\;-\dfrac{1}{3}  t<x<-\dfrac{1}{3} + t\;\Big)\\
    &1 && \Big(\; -\dfrac{1}{3} + t< x < \dfrac{1}{2}t\;\Big)\\
    &0 && \Big(\;  +\dfrac{1}{3} + \dfrac{1}{2}t <x \;\Big)
  \end{aligned}\right. \labeqn{35}
}
\begin{figure}[!htbp]
  \centering
  \includegraphics[width=0.5\textwidth]{./image_tex/beftcrit.png}
  \caption{\refeqn{35}式で表される解。}
  \label{fig:8}
\end{figure}

\vs{0}
　今\refeqn{35}式の膨張派と衝撃波はそれぞれがRiemann問題と見做せるような領域で積分形のBurgers方程式を満たし、ここに積分の加法性\footnote{
  $\int_a^b = \int_a^m + \int_m^b$
}を加味すればその組み合わせとして表される\refeqn{35}式もまた積分形のBurgers方程式を満たすことがわかる。よって、\refeqn{35}式で表される解は弱解である。
\FloatBarrier
\vs{0}
　ただし、\reffig{8}を見ると膨張波は衝撃波よりも早く進むことがわかり、計算すると$t_c = 4/3$で追いついてしまう。（この時衝撃波面は$x_c=1$にある。）この時、膨張波が追いつくと$-1/3 + t < x< 1/2 t$を満たす領域（すなわち膨張波と衝撃波面の間の$u=1$の部分、\reffig{8}で表すと、右二つの赤点の間の「平坦」な部分）が消失してしまい、「系が二つのRiemann問題の足し合わせと見做せる」という前提が崩れてしまう。よって$t>t_c$以降では\refeqn{35}式で表される解が弱解となるとは限らない。
\vs{0}
　そこで、$t>t_c$以降の弱解の候補として\refeqn{36}式及び\reffig{9}で表されるような解を考える。\footnote{
  衝撃波面$x_s(t)$が必ずしも一定でない速度で進むと仮定し、また衝撃波の後方は$t<t_c$の膨張波解と同様の式で表されると仮定して\refeqn{36}式を導いた。
}
\neqn{
  u(x) = \left\{
  \begin{aligned}
    \;&\dfrac{x+1/3}{t}  \qquad && \Big(\; -\dfrac{1}{3}<x<x_s\bracketnormal{t}\;\Big)\\
    &0   && \Big(\;\text{otherwise}\;\Big)
  \end{aligned}\right. \labeqn{36}
}
\begin{figure}[!htbp]
  \centering
  \includegraphics[width=0.7\textwidth]{./image_tex/aftcrit.png}
  \caption{\refeqn{36}式で表される解。}
  \label{fig:9}
\end{figure}

この時、この解が積分形のBurgers方程式
\neqn{
  \difop{t}\dint_{x_1}^{x_2} U dx = - \dfrac{u(x_2,t)^2}{2}+\dfrac{u(x_1,t)^2}{2}\labeqn{37}
}
の解となる条件を考える。以降、見やすさのため膨張波の起点を$x_0=1/3$と書くことにする。（すなわち膨張波は$U = (x-x_0)/t$となる。）この時、\refeqn{37}式は$-1/3 < x_1 < x_s(t)$、$x_s(t)<x_2$として確かめれば十分で、この時\refeqn{37}式の左辺及び右辺は
\neqns{lcl}{
  \text{LHS} &\neqnsequal &\difop{t}\dint_{x_1}^{x_s(t)} U dx+\difop{t}\dint_{x_s(t)}^{x_2} U dx\\
  &\neqnsequal & \difop{t}\dint_{x_1}^{x_s(t)} U dx\\
  &\neqnsequal & \difop{t}\dint_{x_1}^{x_s(t)} \dfrac{x+x_0}{t} dx\\
  &\neqnsequal & \dfrac{x_s(t)-x_0}{t}\dif{x_s(t)}{t} - \dfrac{1}{2}\dfrac{\bracketnormal{x_s(t)-x_0}^2}{t^2} + \dfrac{1}{2}\dfrac{\bracketnormal{x_1-x_0}^2}{t^2}\\[10pt]
  \text{RHS} &\neqnsequal & \dfrac{u(x_2,t)^2}{2}+\dfrac{u(x_1,t)^2}{2}\\
  &\neqnsequal & \dfrac{u(x_1,t)^2}{2}\\
  &\neqnsequal & \dfrac{1}{2}\dfrac{\bracketnormal{x_1-x_0}^2}{t^2}
}
となるので、$\text{LHS}=\text{RHS}$から$x_s(t)$が満たすべき条件が次のようになる。
\neqn{
  \dif{x_s(t)}{t} = \dfrac{x_s(t)-x_0}{2t}
}
よってこの微分方程式を初期条件$x(t_c=4/3) = x_c=1$の下とけば、最終的に$x_s(t)$として次を得る。
\neqn{
  x_s(t) = x_0 + \dfrac{2}{\sqrt{3}}\sqrt{t}\labeqn{46}
}
\FloatBarrier
\vs{0}
よって、$x_s(t)$が\refeqn{46}式で表される時\refeqn{36}式で表される解は$t>t_c$における弱解となる。
\vs{0}
　よって、$t<t_c$及び$t_c$での弱解を組み合わせれば、最終的に弱解として次を得る。
\neqns{l}{
  u(x,t<t_c) = \left\{
  \begin{aligned}
    \;&0  && \Big(\; x<-\dfrac{1}{3}\;\Big)\\
    &\dfrac{x+1/3}{t}  \qquad && \Big(\;-\dfrac{1}{3}  t<x<-\dfrac{1}{3} + t\;\Big)\\
    &1 && \Big(\; -\dfrac{1}{3} + t< x < \dfrac{1}{2}t\;\Big)\\
    &0 && \Big(\;  +\dfrac{1}{3} + \dfrac{1}{2}t <x \;\Big)
  \end{aligned}\right. \labeqn{48}\\[10pt]
  u(x,t>t_c) = \left\{
  \begin{aligned}
    \;&\dfrac{x+1/3}{t}  \qquad && \Big(\; -\dfrac{1}{3}<x<\dfrac{1}{3} + \dfrac{2}{\sqrt{3}}\sqrt{t}\;\Big)\\
    &0   && \Big(\;\text{otherwise}\;\Big)
  \end{aligned}\right. \labeqn{49}
}
また、これを図に表したのが\reffig{10}である。
\begin{figure}[!htbp]
  \centering
  \includegraphics[width=0.9\textwidth]{./image_tex/analytic.png}
  \caption{考えている弱解の模式図。}
  \label{fig:10}
\end{figure}
\vs{0}
　\reffig{10}を見ると、弱解は$t<t_c$においては一定速度進む衝撃波とその後ろからやはり一定速度で進む膨張波の組み合わせとして表され、$t=t_c$で膨張波が衝撃波に追いつくと、以降は衝撃波面が減速しながら高さを失っていき、やがて完全に平坦な解$u=0$に収束することがわかる。
\FloatBarrier

\vs{1}
　さて、弱解がもとまったので、次に考えるべきはこの弱解がエントロピー条件を満たすか否かである。そこで、得られた弱解を参考に特性曲線を描いてみると、\reffig{11}のようになる。
\begin{figure}[!htbp]
  \centering
  \includegraphics[width=0.75\textwidth]{./image_tex/rieman_curve.png}
  \caption{考えている弱解の特性曲線$dx/dt=u(x,t)$。特性曲線に沿って$u$の値が保存する。オレンジの線は衝撃波面の位置を表し、青の領域は膨張波を表す。また緑は衝撃波と膨張波に挟まれた$u=1$の平坦部を表し、これは$t=t_c$で消滅する。また、灰色の領域は$u=0$となっている部分である。}
  \label{fig:11}
\end{figure}
\vs{0}
　\reffig{11}をみると、特性曲線はオレンジで表される衝撃波面に吸収されていることがわかる。よって特性曲線が不連続に収束している（不連続から特性曲線が湧き出していない）ことがわかり、今考えている弱解はエントロピー条件を満たしている。よって、\refeqn{48}\refeqn{49}式（あるいは\reffig{10}）で表される解が今回の系における解析解である。
\FloatBarrier
\subsection*{計算}
　時間ステップ$n$において$u_i^n$が求まっている時、一次精度風上差分法及び2段階Lax-Wendroff法による$u_i^{n+1}$の更新スキームは次のよう。
\neqns{ll}{
  \text{1次精度風上差分:} \qquad & u_i^{n+1}=u_i^n-\frac{\Delta t}{\Delta x}\left(\hat f_{i+1/2}^n-\hat f_{i-1/2}^n\right)\\[5pt]
  &  \qquad \hat f_{i+1/2}^n=\frac{(u_i^n)^2}{2}\\[10pt]
  \text{２段階Lax-Wendroff:}\qquad & u_i^{n+1}=u_i^n-\frac{\Delta t}{\Delta x}\left(\hat f_{i+1/2}^n-\hat f_{i-1/2}^n\right)\\[5pt]
  & \qquad\hat f_{i+1/2}^n=\frac{1}{2}\left(u_{i+1/2}^{n+1/2}\right)^2-\gamma_{i+1/2}^n\frac{\Delta x}{\Delta t}\left(u_{i+1}^n-u_i^n\right)\\[5pt]
  &\qquad u_{i+1/2}^{n+1/2}=\frac{u_i^n+u_{i+1}^n}{2}-\frac{\Delta t}{4\Delta x}\left[(u_{i+1}^n)^2-(u_i^n)^2\right]\\[5pt]
  &\qquad\gamma_{i+1/2}^n=\epsilon\left|u_{i+1}^n-u_i^n\right|
}
　特に2段階Lax-Wendroffにおいては$u$の勾配が急なところで強くなるような人工粘性$\gamma$を導入している。すなわち$\gamma = \epsilon \cdot \partial U/\partial x$として人工粘性を決定する。ただし$\epsilon$は人工粘性の強さを決めるパラメタである。
\vs{0}
　また、計算に当たっては$\Delta t = \nu \Delta x$によって$\Delta t$を定めることにする。（$\nu$は適切に決定する。）また、$\Delta x =\bracketnormal{\text{領域幅}}/\bracketnormal{\text{格子数}}$として$\Delta x$を定めることにする。すなわち空間方向の格子数$M$が精度を決定するパラメタとなる。
\vs{0}
　以上より、今回の計算においてパラメタとなるのは格子数$M$と人工粘性の強さ$\epsilon$の２種類である。よってこれらのパラメタを適宜調整しながら結果を観察する。
\subsection*{結果}
　まず、\reffig{100}に1次精度風上差分法を用いた計算結果を示す。（\href{https://shota-mitamura.com/works/geo_analysiskadai1/kadai/}{動画ページ}では図７に対応している。）\reffig{100}の一番左の図は格子数を$50$として計算した場合の$t=1$頃の結果を示しているが、この画像を見ると衝撃波及び膨張波は捉えられているものの、数値拡散により精度が悪く衝撃波面の傾きが緩くなってしまっていることがわかる。ただし、格子数を$100,200$としてメッシュを細かくしていくと数値拡散は弱くなっていき、数値解の精度が良くなっていく。（\reffig{100}中央図及び右図参照。）
\vs{0}
　次に、\reffig{101}に2段階Lax-Wendroff法を用いた計算結果を示す。（\href{https://shota-mitamura.com/works/geo_analysiskadai1/kadai/}{動画ページ}では図８に対応している。）ただし格子数は$100$で固定しており、人工粘性の強さ$\epsilon$を変化させている。まず\reffig{101}左図の一番人工粘性が強い場合（$\epsilon =0.05$）を観察すると、一次精度風上差分に比べて数値拡散が抑えられており、一次精度風上差分で同じメッシュ数で計算した場合（\reffig{100}中央図）と比べても衝撃波の急峻な立ち上がりを捉えることができていることがわかる。ただしその対価として数値振動がが発生してしまっている。また、\reffig{101}中央図及び右図（それぞれ$\epsilon = 0.2,0.4$、右に行くほど人工粘性が強くなっている）を見ると、人工粘性が強いほど数値振動が抑えられる代わりに数値拡散が発生し衝撃波の立ち上がりを捉えづらくなっていることがわかる。
\begin{figure}[!htbp]
    \centering
    \begin{subfigure}{0.31\textwidth}
        \centering
        \includegraphics[width = \linewidth]{./image_tex/burgersupwind3000032.png}
    \end{subfigure}
    \begin{subfigure}{0.31\textwidth}
        \centering
        \includegraphics[width=\linewidth]{./image_tex/burgersupwind2000063.png}
    \end{subfigure}
    \begin{subfigure}{0.31\textwidth}
        \centering
        \includegraphics[width=\linewidth]{./image_tex/burgersupwind1000125.png}
    \end{subfigure}
    \caption{1次精度風上差分を用いた数値計算の結果。画像の右にいくほど空間の格子数が多くなっている。（すなわち$\Delta x$が小さくなっている。）}
    \label{fig:100}
\end{figure}

\begin{figure}[!htbp]
    \centering
    \begin{subfigure}{0.31\textwidth}
        \centering
        \includegraphics[width = \linewidth]{./image_tex/burgerslax4000063.png}
    \end{subfigure}
    \begin{subfigure}{0.31\textwidth}
        \centering
        \includegraphics[width=\linewidth]{./image_tex/burgerslax6000063.png}
    \end{subfigure}
    \begin{subfigure}{0.31\textwidth}
        \centering
        \includegraphics[width=\linewidth]{./image_tex/burgerslax5000063.png}
    \end{subfigure}
    \caption{2段階Lax-Wendroffを用いた数値計算の結果。格子数は固定し、人工粘性の強さ$\epsilon$を変化させている。画像の右にいくほど人工粘性が強くなっている。}
    \label{fig:101}
\end{figure}
\FloatBarrier
\subsection*{スキームの比較}
　最後に二つのスキームの比較を行う。\reffig{102}は1次精度風上差分（青）及び2段階Lax-Wendroff（オレンジ）による数値解を同じ画面に表示したものであり、左図は膨張波が衝撃波に追いつく前$\bracketnormal{t<t_c}$での様子で右図は膨張波が衝撃波に追いついた後$\bracketnormal{t>t_c}$での様子ある。（\href{https://shota-mitamura.com/works/geo_analysiskadai1/kadai/}{動画ページ}では図９に対応している。）この図や動画を見ると、いずれのスキームも衝撃波の速度を正しく捉えることができており、それは衝撃波の速度が一定で無くなった$t>t_c=4/3$以降も変わらない。（ただし精度の悪さに由来するずれは生じている。）また、やはり１次元風上差分では数値拡散が生じる影響で衝撃波面の立ち上がりが鈍くなっており、Lax-Wendroffの方が衝撃波面がより解析解に近くなっている。
\vs{0}
　また、メッシュを細かくするといずれのスキームでも衝撃波の再現精度が良くなることがわかる。（\reffig{103}参照。\href{https://shota-mitamura.com/works/geo_analysiskadai1/kadai/}{動画ページ}では図１０に対応している。）
　
\begin{figure}[!htbp]
    \centering
    \begin{subfigure}{0.45\textwidth}
        \centering
        \includegraphics[width = \linewidth]{./image_tex/burgerscombined000063.png}
    \end{subfigure}
    \hspace{0.05\linewidth}
    \begin{subfigure}{0.45\textwidth}
        \centering
        \includegraphics[width=\linewidth]{./image_tex/burgerscombined000175.png}
    \end{subfigure}
    \caption{1次精度風上差分（青）及び2段階Lax-Wendroff（オレンジ）の比較。左は膨張波が衝撃波に追いつく前での様子。右は膨張波が衝撃波に追いついた後の様子。}
    \label{fig:102}
\end{figure}
\begin{figure}[!htbp]
    \centering
    \begin{subfigure}{0.45\textwidth}
        \centering
        \includegraphics[width = \linewidth]{./image_tex/burgerscombined2000125.png}
    \end{subfigure}
    \hspace{0.05\linewidth}
    \begin{subfigure}{0.45\textwidth}
        \centering
        \includegraphics[width=\linewidth]{./image_tex/burgerscombined2000350.png}
    \end{subfigure}
    \caption{\reffig{102}よりもメッシュを細かくした場合の数値計算結果。}
    \label{fig:103}
\end{figure}
\FloatBarrier


\newpage
\section{参考文献}
\begin{thebibliography}{9}

\bibitem{Kawakuwa}
河村・桑原スキーム

\url{https://pbcglab.jp/cgi-bin/wiki/?\%E6\%B2\%B3\%E6\%9D\%91\%E3\%83\%BB\%E6\%A1\%91\%E5\%8E\%9F\%E3\%82\%B9\%E3\%82\%AD\%E3\%83\%BC\%E3\%83\%A0}

\bibitem{QUICKEST}
QUICKESTスキーム

\url{https://hydroeurope.upc.edu/wp-content/uploads/2025/04/QUICKEST-Numerical-scheme.pdf}

\bibitem{QUICKESTorder}
QUICKESTスキームの空間精度についての参照記事

\url{https://yoshida-geophys.jp/comp/adscheme.pdf}

\end{thebibliography}

\end{document}


